% =====================================================================
%  Shared preamble for Math F113X solution sets.
%
%  One preamble drives all three documents. The body file selects its
%  flavor with a single command right after \begin{document} -- or by
%  the build script setting the switches for it:
%
%     \reviewtrue      markup + index page      (EXPANDED, review copy)
%     \reviewfalse     markup stripped          (EXPANDED, student copy)
%     \transcriptstyle green boxes, no markup   (TRANSCRIPT)
%
%  Environments:
%     stmt   problem statement          (grey box, both documents)
%     sol    the solution               (blue, or green in transcript)
%     add    an addition to the         (blue rule + [+n], indexed)
%            handwritten solution
%     corr   a correction to a          (orange rule + [fix n], indexed)
%            handwritten error
%     flag   an error left in place     (orange box, transcript only)
%     \nohw  marks a problem with no handwritten solution
% =====================================================================

\documentclass[11pt, oneside]{article}

\usepackage{geometry}
\geometry{letterpaper, margin=1in}
\usepackage[parfill]{parskip}
\usepackage{graphicx}
\usepackage{amssymb, amsmath}
\usepackage{array}
\usepackage{enumitem}
\usepackage[dvipsnames]{xcolor}
\usepackage[most]{tcolorbox}
\usepackage{needspace}

\newif\ifreview
\reviewtrue          % the build script flips this for the clean copy

% ---- problem header --------------------------------------------------
\newcommand{\problem}[1]{%
  \vspace{1.6em}
  \noindent\rule{\textwidth}{0.8pt}\\[0.3em]
  \noindent{\bf\large Problem #1}\\[0.3em]
  \noindent\rule{\textwidth}{0.4pt}
  \vspace{0.4em}
}

% ---- part label; Needspace keeps it with the start of its solution ----
\newcommand{\prt}[1]{\Needspace*{6\baselineskip}\vspace{1.1em}%
  \noindent{\bf\large (#1)}\ }

% ---- problem statement -----------------------------------------------
\newtcolorbox{stmt}{%
  enhanced, breakable,
  colback=black!2, colframe=black!25,
  boxrule=0.5pt, arc=2mm,
  left=3.5mm, right=3.5mm, top=2.5mm, bottom=2.5mm,
  before upper={\setlength{\parskip}{0.5\baselineskip}},
  before skip=8pt, after skip=8pt
}

% ---- the solution ----------------------------------------------------
\definecolor{solframe}{named}{RoyalBlue}
\newcommand{\soltitle}{Solution}

\newtcolorbox{sol}[1][\soltitle]{%
  enhanced, breakable,
  colback=solframe!4, colframe=solframe!55!black,
  boxrule=0.7pt, arc=3.5mm,
  left=4mm, right=4mm, bottom=3mm, top=4.5mm,
  coltitle=white, colbacktitle=solframe!55!black,
  fonttitle=\bfseries\small,
  attach boxed title to top left={xshift=5mm, yshift=-2.6mm},
  boxed title style={arc=1.8mm, boxrule=0pt, left=2.5mm, right=2.5mm},
  title={#1},
  before upper={\setlength{\parskip}{0.55\baselineskip}},
  before skip=10pt, after skip=10pt
}

% Transcript documents call this once, before \begin{document}.
% Green signals "this is the handwritten work, untouched."
\newcommand{\transcriptstyle}{%
  \colorlet{solframe}{ForestGreen}%
  \renewcommand{\soltitle}{Handwritten solution}%
  \reviewfalse
}

% ---- indexed additions and corrections -------------------------------
\newcounter{addno}
\newcounter{fixno}
\newcommand{\additem}[2]{\noindent\makebox[3.2em][l]{%
  {\sffamily\bfseries\color{RoyalBlue!65!black}[+#1]}}\parbox[t]{0.85\textwidth}{#2}\par\vspace{0.35em}}
\newcommand{\fixitem}[2]{\noindent\makebox[3.2em][l]{%
  {\sffamily\bfseries\color{BurntOrange!65!black}[fix #1]}}\parbox[t]{0.85\textwidth}{#2}\par\vspace{0.35em}}
% The index is written to its own file as the document is typeset, then closed
% and read back on the index page. Because every add/corr occurs before the
% index page, this resolves within a single pdflatex pass.
\newwrite\adlfile
\newwrite\fxlfile
\AtBeginDocument{%
  \immediate\openout\adlfile=\jobname.adl
  \immediate\openout\fxlfile=\jobname.fxl}

\newcommand{\adlist}{%
  \immediate\closeout\adlfile
  \ifnum\value{addno}=0 \emph{None.}\else\input{\jobname.adl}\fi}
\newcommand{\fxlist}{%
  \immediate\closeout\fxlfile
  \ifnum\value{fixno}=0
    \emph{None --- no errors were found in the handwritten pages for this assignment.}%
  \else\input{\jobname.fxl}\fi}

% #1 = the one-line description that appears in the index
\newenvironment{add}[1]
{\par\ifreview
   \refstepcounter{addno}%
   \immediate\write\adlfile{\string\additem{\theaddno}{\unexpanded{#1}}}%
   \begin{tcolorbox}[blanker, breakable,
     borderline west={2.2pt}{0pt}{RoyalBlue!35},
     left=5mm, right=0mm, top=1.2mm, bottom=1.2mm,
     before skip=7pt, after skip=7pt,
     before upper={\setlength{\parskip}{0.5\baselineskip}}]%
   \noindent{\scriptsize\sffamily\color{RoyalBlue!65!black}[+\theaddno]}\ %
 \fi\ignorespaces}
{\ifreview\end{tcolorbox}\fi}

\newenvironment{corr}[1]
{\par\ifreview
   \refstepcounter{fixno}%
   \immediate\write\fxlfile{\string\fixitem{\thefixno}{\unexpanded{#1}}}%
   \begin{tcolorbox}[blanker, breakable,
     borderline west={2.2pt}{0pt}{BurntOrange!60},
     left=5mm, right=0mm, top=1.2mm, bottom=1.2mm,
     before skip=7pt, after skip=7pt,
     before upper={\setlength{\parskip}{0.5\baselineskip}}]%
   \noindent{\scriptsize\sffamily\color{BurntOrange!65!black}[fix \thefixno]}\ %
 \fi\ignorespaces}
{\ifreview\end{tcolorbox}\fi}

% ---- transcript-only: an error left in place and flagged --------------
\newtcolorbox{flag}[1][Transcription note --- possible error]{%
  enhanced, breakable,
  colback=BurntOrange!7, colframe=BurntOrange!75!black,
  boxrule=0.7pt, arc=3.5mm,
  left=4mm, right=4mm, top=4.5mm, bottom=3mm,
  fonttitle=\bfseries\small,
  coltitle=white, colbacktitle=BurntOrange!75!black,
  attach boxed title to top left={xshift=5mm, yshift=-2.6mm},
  boxed title style={arc=1.8mm, boxrule=0pt, left=2.5mm, right=2.5mm},
  title={#1},
  before upper={\setlength{\parskip}{0.55\baselineskip}},
  before skip=10pt, after skip=10pt
}
\newcommand{\fix}[1]{\vspace{0.3em}\noindent{\bf Suggested fix:} #1}

% ---- misc ------------------------------------------------------------
\newcommand{\nohw}{\ifreview{\scriptsize\sffamily\color{RoyalBlue!65!black}%
  [entirely new --- no handwritten solution exists for this problem]}%
  \par\vspace{0.3em}\fi}
\newcommand{\nosol}{\vspace{0.4em}\noindent{\color{Gray}\itshape
  [No handwritten solution for this problem in the folder.]}\vspace{0.4em}}
\newcommand{\ans}[1]{\fbox{#1}}
\newcommand{\crit}[1]{\underline{#1}}   % critical player in a coalition
\newcommand{\srcnote}[1]{\vspace{0.2em}\noindent{\small\color{Gray}\itshape
  Source: #1}\vspace{0.2em}}

\reviewfalse
\usepackage{tikz}
\usepackage{multicol}
\usetikzlibrary{calc}

\tikzset{
  vtx/.style={circle, draw, fill=white, inner sep=1.5pt, minimum width=9pt, font=\scriptsize},
  wl/.style={midway, inner sep=1.5pt, fill=white, font=\scriptsize},
  wlb/.style={midway, inner sep=1.5pt, fill=white, font=\small},
}

% ---- a Nearest-Neighbour route drawn as a chain, the way the pages draw it.
%      #1 = "F/19, A/43, M/38, P/41, D/"   (last weight empty)
%      #2 = weight of the closing edge, drawn as an arc underneath
\newcounter{rtstep}
\newcommand{\route}[2]{%
\setcounter{rtstep}{0}%
\begin{tikzpicture}[baseline=-2pt, every node/.style={font=\small}]
\foreach \v/\w [count=\i from 0] in {#1}{%
  \stepcounter{rtstep}%
  \pgfmathsetmacro\xx{1.62*\i}%
  \ifx\w\empty\else
    \draw (\xx+0.34,0) -- (\xx+1.28,0);
    \node[above, inner sep=1.5pt, font=\footnotesize] at (\xx+0.81,0.03) {\w};
  \fi
  \node[draw, circle, inner sep=1pt, minimum width=20pt, fill=white] at (\xx,0) {\v};}
\pgfmathsetmacro\lx{1.62*(\value{rtstep}-1)}
\draw (0,-0.32) .. controls (0.35,-1.18) and (\lx-0.35,-1.18) .. (\lx,-0.32);
\node[below, inner sep=1.5pt, font=\footnotesize] at (\lx*0.5,-0.96) {#2};
\end{tikzpicture}}

% ---- #19: the five Dallas-area stores. #1 = highlight colour, #2 = edges
\newcommand{\Dallas}[2][RoyalBlue!22]{%
\begin{tikzpicture}[scale=1.05, every node/.style={font=\small}]
\coordinate (FW) at (-2.6,0.8);
\coordinate (P)  at (0,2.7);
\coordinate (M)  at (2.6,0.8);
\coordinate (A)  at (1.6,-2.2);
\coordinate (D)  at (-1.6,-2.2);
\foreach \i/\j in {#2}{\draw[#1, line width=7pt, line cap=round] (\i) -- (\j);}
\draw (FW)--(P) (P)--(M) (M)--(A) (A)--(D) (D)--(FW);
\draw (FW)--(M) (FW)--(A) (P)--(A) (P)--(D) (M)--(D);
\node at (-1.45,1.90) {54};  \node at ( 1.45,1.90) {38};
\node at ( 2.30,-0.70) {43}; \node at ( 0.00,-2.45) {50};
\node at (-2.30,-0.70) {42}; \node at ( 0.00, 1.05) {52};
\node at (-0.55,-0.90) {19}; \node at ( 0.95, 0.30) {53};
\node at (-0.95, 0.30) {41}; \node at ( 0.60,-0.35) {56};
\tikzstyle{v}=[draw, circle, fill=black, inner sep=1.7pt]
\node[v,label=left:{Fort Worth}]  at (FW) {};
\node[v,label=above:{Plano}]      at (P)  {};
\node[v,label=right:{Mesquite}]   at (M)  {};
\node[v,label=right:{Arlington}]  at (A)  {};
\node[v,label=below:{Denton}]     at (D)  {};
\end{tikzpicture}}

% ---- #22 Sorted Edges: only the seven chosen edges, laid out as the page draws them
\newcommand{\Israel}{%
\begin{tikzpicture}[scale=1.5, every node/.style={font=\small}]
\coordinate (N) at (0.55,2.05); \coordinate (T) at (1.95,2.35);
\coordinate (H) at (0,0.45);    \coordinate (TA) at (4.35,0.75);
\coordinate (J) at (3.55,2.30); \coordinate (E) at (1.45,-0.85);
\coordinate (BS) at (3.45,-0.85);
\draw (N)--node[wlb]{29}(T);
\draw (H)--node[wlb]{35}(N);
\draw (H)--node[wlb,pos=0.35]{95}(TA);
\draw (J)--node[wlb]{58}(TA);
\draw (J)--node[wlb,pos=0.55]{81}(BS);
\draw (T)--node[wlb,pos=0.45]{405}(E);
\draw (E)--node[wlb]{241}(BS);
\tikzstyle{v}=[draw, circle, fill=black, inner sep=1.7pt]
\node[v,label=above left:{N}] at (N){};
\node[v,label=above:{T}]       at (T){};
\node[v,label=left:{H}]        at (H){};
\node[v,label=right:{TA}]      at (TA){};
\node[v,label=above:{J}]       at (J){};
\node[v,label=left:{E}]        at (E){};
\node[v,label=right:{BS}]      at (BS){};
\end{tikzpicture}}

% ---- Problem A: the pentagon. #1 = highlight colour, #2 = circuit edges
\newcommand{\PentA}[2][RoyalBlue!25]{%
\begin{tikzpicture}[scale=0.92]
\node[vtx,minimum width=15pt,font=\small] (B) at (90:3){$B$};
\node[vtx,minimum width=15pt,font=\small] (C) at (18:3){$C$};
\node[vtx,minimum width=15pt,font=\small] (E) at (306:3){$E$};
\node[vtx,minimum width=15pt,font=\small] (D) at (234:3){$D$};
\node[vtx,minimum width=15pt,font=\small] (A) at (162:3){$A$};
\foreach \i/\j in {#2}{\draw[#1, line width=7pt, line cap=round] (\i) -- (\j);}
\foreach \i/\j/\k/\p in {A/B/1/0.5, A/C/1/0.5, B/C/1/0.5, D/E/100/0.5,
                         A/D/2/0.5, A/E/2/0.32, B/D/2/0.32, B/E/2/0.32,
                         C/D/2/0.32, C/E/2/0.5}
  {\draw (\i) --node[pos=\p, inner sep=1.5pt, fill=white, font=\scriptsize]{\k} (\j);}
\foreach \n in {A,B,C,D,E}{\node[vtx,minimum width=15pt,font=\small] at (\n){$\n$};}
\end{tikzpicture}}

\title{Math F113X: Homework Set 6b --- Solutions}
\author{}
\date{}

\begin{document}
\maketitle
\vspace{-3em}

\begin{center}
\fbox{\parbox{0.9\textwidth}{\centering
Graph Theory: \# 19*, 22*, A, B, C\\[0.4em]
{\small * In part b of both \#19 and \#22 you were asked to run Nearest Neighbor
from only \textbf{two} additional vertices. All of them are worked below.}}}
\end{center}


%%%%%%%%%%%%%%%%%%%%%%%%%%%%%%%%%%%%%%%%%%%%%%%%%%%%%%%%%%%%
\problem{19}

\begin{stmt}
A company needs to deliver product to each of their 5 stores around the Dallas, TX area. Driving
distances between the stores are shown below. Find a route for the driver to follow, returning to
the distribution center in Fort Worth:

\begin{enumerate}[label=\alph*., nosep]
\item Using Nearest Neighbor starting in Fort Worth
\item Using Repeated Nearest Neighbor
\item Using Sorted Edges
\end{enumerate}

\begin{center}
\begin{tabular}{|l|c|c|c|c|}
\hline
& \textbf{Plano} & \textbf{Mesquite} & \textbf{Arlington} & \textbf{Denton} \\
\hline
\textbf{Fort Worth} & 54 & 52 & 19 & 42 \\
\textbf{Plano}      &    & 38 & 53 & 41 \\
\textbf{Mesquite}   &    &    & 43 & 56 \\
\textbf{Arlington}  &    &    &    & 50 \\
\hline
\end{tabular}
\end{center}
\end{stmt}

\prt{a}
\begin{sol}
\begin{center}\route{F/19, A/43, M/38, P/41, D/}{42}\end{center}

\vspace{0.3em}
Weight: \ans{183}
\end{sol}

\prt{b}
\begin{sol}
Starting vertex:

\begin{center}
\begin{tabular}{@{}l@{\hspace{1.5em}}l@{}}
\route{F/19, A/43, M/38, P/41, D/}{42} & wgt: 183 \ (part a) \\[1.9em]
\route{A/19, F/42, D/41, P/38, M/}{43} & wgt: 183 \\[1.9em]
\route{P/38, M/43, A/19, F/42, D/}{41} & wgt: 183 \\[1.9em]
\route{M/38, P/41, D/42, F/19, A/}{43} & wgt: 183 \\[1.9em]
\route{D/41, P/38, M/43, A/19, F/}{42} & wgt: 183 \\
\end{tabular}
\end{center}

\vspace{0.5em}
\underline{ANS}: Pick $F$--$A$--$M$--$P$--$D$--$F$, \ wgt \ans{183}.

\begin{add}{One line saying why every starting vertex lands on the same number, which the table of five identical answers shows but does not say.}
Every start produces the same circuit, just entered at a different point, so all five weigh 183.
\end{add}
\end{sol}

\clearpage
\prt{c}
\begin{sol}
Sorted Edges:

\begin{center}
\begin{tabular}{c@{\hspace{3em}}c}
\Dallas[ForestGreen!25]{FW/A, A/M, M/P, P/D, D/FW} &
$\begin{array}{r}
19 \\ 38 \\ 41 \\ 42 \\ 43 \\ \hline 183
\end{array}$ \\
\end{tabular}
\end{center}

\begin{corr}{\#19(c): the page lists the fifth edge as 53 (Plano--Arlington) and draws it, but 53 does not fit the total of 183 written beside it. Sorted Edges takes Mesquite--Arlington at 43: after 19, 38, 41 and 42 the chain is Arlington--Fort Worth--Denton--Plano--Mesquite, and 43 is the cheapest edge that closes it. $19+38+41+42+43 = 183$, which is the total the page already has.}
The edges are taken in order of weight: 19 ($FW$--$A$), 38 ($P$--$M$), 41 ($P$--$D$),
42 ($FW$--$D$), and finally 43 ($M$--$A$), which closes the circuit.
\end{corr}

\begin{add}{Notes that this is the same circuit the first two parts found, which is the point of the problem.}
This is the same circuit as in parts (a) and (b), so all three algorithms agree here.
\end{add}
\end{sol}

%%%%%%%%%%%%%%%%%%%%%%%%%%%%%%%%%%%%%%%%%%%%%%%%%%%%%%%%%%%%
\problem{22}

\begin{stmt}
A tourist wants to visit 7 cities in Israel. Driving distances, in kilometers, between the cities
are shown below. Find a route for the person to follow, returning to the starting city:

\begin{enumerate}[label=\alph*., nosep]
\item Using Nearest Neighbor starting in Jerusalem
\item Using Repeated Nearest Neighbor
\item Using Sorted Edges
\end{enumerate}

\begin{center}\small
\begin{tabular}{|l|c|c|c|c|c|c|}
\hline
& \textbf{Jerusalem} & \textbf{Tel Aviv} & \textbf{Haifa} & \textbf{Tiberias} & \textbf{Beer Sheba} & \textbf{Eilat} \\
\hline
\textbf{Jerusalem}  & --  & 58  & 151 & 152 & 81  & 309 \\
\textbf{Tel Aviv}   & 58  & --  & 95  & 134 & 105 & 346 \\
\textbf{Haifa}      & 151 & 95  & --  & 69  & 197 & 438 \\
\textbf{Tiberias}   & 152 & 134 & 69  & --  & 233 & 405 \\
\textbf{Beer Sheba} & 81  & 105 & 197 & 233 & --  & 241 \\
\textbf{Eilat}      & 309 & 346 & 438 & 405 & 241 & --  \\
\textbf{Nazareth}   & 131 & 102 & 35  & 29  & 207 & 488 \\
\hline
\end{tabular}
\end{center}
\end{stmt}

\prt{a}
\begin{sol}
\begin{center}\route{J/58, TA/95, H/35, N/29, T/233, BS/241, E/}{309}\end{center}

\vspace{0.3em}
The six edges along the top come to 691.

\begin{corr}{\#22(a): the page stops at 691, which is the weight of the \emph{path} $J$--$TA$--$H$--$N$--$T$--$BS$--$E$. The tour has to return to Jerusalem, and that leg costs 309. The next page does add it --- the Jerusalem row of part (b) reads 1000 --- so this is an unfinished line rather than a wrong one.}
Adding the return leg $E$--$J = 309$ gives a circuit of weight \ans{1000}.
\end{corr}
\end{sol}

\prt{b}
\begin{sol}
Starting vertex, and weight:

\begin{center}
\begin{tabular}{@{}l@{\hspace{1.5em}}r@{}}
\route{J/58, TA/95, H/35, N/29, T/233, BS/241, E/}{309}     & 1000 \\[1.9em]
\route{TA/58, J/81, BS/197, H/35, N/29, T/405, E/}{346}     & 1151 \\[1.9em]
\route{H/35, N/29, T/134, TA/58, J/81, BS/241, E/}{438}     & 1016 \\[1.9em]
\route{T/29, N/35, H/95, TA/58, J/81, BS/241, E/}{405}      & 944 \\[1.9em]
\route{BS/81, J/58, TA/95, H/35, N/29, T/405, E/}{241}      & 944 \\[1.9em]
\route{E/241, BS/81, J/58, TA/95, H/35, N/29, T/}{405}      & 944 \\[1.9em]
\route{N/29, T/69, H/95, TA/58, J/81, BS/241, E/}{488}      & 1061 \\
\end{tabular}
\end{center}

\vspace{0.5em}
Pick \ $T$--$N$--$H$--$TA$--$J$--$BS$--$E$, \ wght: \ans{944}

\begin{add}{Names the three starts that tie, since the table shows them but the ``Pick'' line reports only one.}
Tiberias, Beer Sheba, and Eilat all give a circuit with weight 944 because they produce the same
circuit entered at three different starting points.
\end{add}
\end{sol}

\prt{c}
\begin{sol}
\begin{center}
\begin{tabular}{c@{\hspace{2.5em}}c}
\Israel &
$\begin{array}{r}
29 \\ 35 \\ 95 \\ 58 \\ 81 \\ 405 \\ 241 \\ \hline 944
\end{array}$ \\
\end{tabular}
\end{center}

\begin{add}{Records that Sorted Edges lands on the same circuit as the best Repeated Nearest Neighbor run, which is what makes 944 the answer to report.}
This is the same circuit that part (b) found from Tiberias, Beer Sheba and Eilat.
\end{add}
\end{sol}

%%%%%%%%%%%%%%%%%%%%%%%%%%%%%%%%%%%%%%%%%%%%%%%%%%%%%%%%%%%%
\clearpage
\problem{A}

\begin{stmt}
Answer questions about the graph below.

\begin{center}\PentA{}\end{center}

\begin{enumerate}[nosep]
\item Apply the Repeated Nearest Neighbor Algorithm and determine the weight of the resulting
circuit. (Note that you are allowed to use symmetry in the graph to make this faster!)
\item Apply the Sorted Edges Algorithm and determine the weight of the resulting circuit.
\item Using your own judgment, find a minimum weight Hamiltonian circuit.
\item Does either algorithm --- Repeated Nearest Neighbor or Sorted Edges --- return a minimum
weight Hamiltonian circuit?
\end{enumerate}
\end{stmt}

\prt{1}
\begin{sol}
Using symmetry, we see that RNN gives the same weight starting at $A$, $B$, or $C$, and $D$ and $E$
will be the same:

\begin{center}
\route{A/1, B/1, C/2, D/100, E/}{2} \quad $= 106$
\end{center}

\begin{center}
\route{D/2, A/1, B/1, C/2, E/}{100} \quad $= 106$
\end{center}
\end{sol}

\prt{2}
\begin{sol}
\begin{center}\PentA[BurntOrange!35]{A/B, B/C, C/E, E/D, D/A}\end{center}

\vspace{0.3em}
Weight \ans{106}

\begin{add}{Records the order the edges come in, since the figure shows only the finished circuit.}
Sorted Edges takes $AB$ and $BC$ at weight 1 and then must skip $AC$, which would close a triangle.
Of the weight-2 edges it takes $AD$ and $CE$, and the circuit can only be closed with
$DE = 100$.
\end{add}
\end{sol}

\prt{3}
\begin{sol}
\begin{center}\PentA[ForestGreen!30]{A/B, B/D, D/C, C/E, E/A}\end{center}

\vspace{0.3em}
Weight $= \ans{9}$

\begin{add}{Says what makes 9 the best possible, which the figure asserts without argument.}
Any circuit has to give $D$ two edges and $E$ two edges. Avoiding the weight-100 edge between them
means each takes two weight-2 edges, for 8, and the fifth edge can then be one of the weight-1
edges --- so 9 cannot be beaten.
\end{add}
\end{sol}

\prt{4}
\begin{sol}
\ans{No.} RNN and SE \underline{do not} give an optimal solution.

\begin{add}{Names the trap in one line, since this is the point the problem is built around.}
Both algorithms take the cheap edges among $A$, $B$ and $C$ first and are then forced to pay the
100 to get from $D$ to $E$; the optimal circuit gives that up front by keeping $D$ and $E$ apart.
\end{add}
\end{sol}

%%%%%%%%%%%%%%%%%%%%%%%%%%%%%%%%%%%%%%%%%%%%%%%%%%%%%%%%%%%%
\problem{B}

\begin{stmt}
Suppose someone wants to visit the capital city of every state in the contiguous 48 states and
Washington DC. So, they will visit 49 cities in total. [NOTE: you will need to use a computational
tool, like WolframAlpha, to complete this problem.]

\begin{enumerate}[nosep]
\item Suppose they want to start and end the 49-city-tour at the same place. How many different
tours are possible? (Your answer should be in \emph{both} factorial notation and in scientific
notation.)
\item Suppose you have a computer that can calculate the length of 1000 49-city-tours in one
second. How long would it take the computer to calculate the length of all possible
49-city-tours? Give your answer in \emph{years}.
\item What does your answer in part 2 indicate about the Brute Force algorithm?
\end{enumerate}
\end{stmt}

\prt{1}
\begin{sol}
$(48)! = \ans{$1.24 \times 10^{61}$}$

\begin{add}{One line on why the exponent is 48 and not 49, which is the only place this part can go wrong.}
Fixing the starting city leaves the other 48 to be ordered, and each tour is counted once this way.
\end{add}
\end{sol}

\prt{2}
\begin{sol}
$$\frac{1000}{\text{sec}} \cdot \frac{60\ \text{sec}}{\text{min}} \cdot
\frac{60\ \text{min}}{1\ \text{hr}} \cdot \frac{24\ \text{hr}}{1\ \text{day}} \cdot
\frac{365\ \text{days}}{\text{year}} = 3.1536 \times 10^{10} \ \text{per year}$$

So
$$\frac{(48!)}{1000 \cdot 60 \cdot 60 \cdot 24 \cdot 365}
= \ans{$3.93 \times 10^{50}$ years} \ !!$$
\end{sol}

\prt{3}
\begin{sol}
Even with a computer, Brute Force will take too long.
\end{sol}

%%%%%%%%%%%%%%%%%%%%%%%%%%%%%%%%%%%%%%%%%%%%%%%%%%%%%%%%%%%%
\problem{C}

\begin{stmt}
A company plans to run a bus continually throughout the day in a loop around 5 tourist stops in a
city, denoted A, B, C, D, and E. The company would like to find a route that minimizes the time to
complete a loop. Driving times in minutes are shown in the table.

\begin{center}
\begin{tabular}{|c||c|c|c|c|c|}
\multicolumn{6}{c}{Driving Times in Minutes}\\
\hline
& A & B & C & D & E \\
\hline\hline
A & -- & 4 & 7 & 3 & 2 \\
\hline
B & 4 & -- & 5 & 6 & 9 \\
\hline
C & 7 & 5 & -- & 8 & 6 \\
\hline
D & 3 & 6 & 8 & -- & 5 \\
\hline
E & 2 & 9 & 6 & 5 & -- \\
\hline
\end{tabular}
\end{center}

\textbf{Using terminology from Graph Theory}, what is the company looking for?
\end{stmt}

\begin{sol}
\nohw
The table gives a driving time for every pair of stops, so it describes a \ans{complete weighted
graph} on 5 vertices: each tourist stop is a vertex, each edge is the trip between two stops, and
the weight of an edge is the driving time in minutes.

The bus has to call at every stop and come back to where it started, which is a
\ans{Hamiltonian circuit}, and the company wants the one that takes the least total time. So they
are looking for a

\begin{center}
\ans{minimum weight Hamiltonian circuit} \ in this graph.
\end{center}

\vspace{0.3em}
\emph{(This is the Traveling Salesman Problem.)}
\end{sol}


\end{document}
